\documentclass[12pt]{article}
\usepackage[utf8]{inputenc}
\usepackage[T1]{fontenc}
\usepackage{hyperref}
\title{Yiddish Sign Language Language Resources for K-12 Curriculum Development}
\author{Compiled by Claude Sonnet 5.5 (claude-sonnet-5-5) for Andr\'as Kornai}
\date{}
\begin{document}
\maketitle

\section{Introduction}
Search results indicate that Yiddish Sign Language (code yds) is probably not an attested language. The ISO 639-3 code was retired as non-existent, effective 2015-01-12 (IANA registry). The entry appears to trace to a footnote in the preface of Oliver Sacks's \emph{Seeing Voices} (1990). Sources report that sign-language experts and Deaf Chassidic community sources know of no such language, and that some families develop home signs. Some older or less careful listings still include it among the sign languages of Israel. A curriculum cannot responsibly be planned around this language on the evidence found.

\section{Primary Educational Resources}
None found.

\section{Linguistic and Reference Materials}
\begin{itemize}
\item Jewish Languages Project, Yiddish Sign Language (critical account of the evidence): \url{https://jewishlanguages.org/yiddish-sign-language}
\item IANA language subtag template for yds: \url{https://www.iana.org/assignments/lang-subtags-templates/yds.txt}
\end{itemize}

\section{Religious and Community Resources}
\begin{itemize}
\item Jewish Deaf Community Center, Deaf Chassidim Community (relevant community context, not a Yiddish Sign Language resource): \url{https://jdcc.org/?p=1510}
\end{itemize}

\section{Digital and Audio Resources}
None found.

\section{Curriculum Development Context}
There is no verified Yiddish Sign Language to teach in. Deaf Jewish learners in Israel, the United States and elsewhere would be served by the established national sign languages and by Yiddish or Hebrew print literacy, which belong to other dossiers. Any further work would first require confirming that a community using this language exists.
\end{document}